% A paper for an Oxford University Press journal, on OUP's own class,
% oup-authoring-template.
%
% `contemporary` is one of the class's page designs; `modern` and
% `traditional` are the others, and the journal's instructions say which.
% `webpdf` makes the links work in the PDF.

\documentclass[unnumsec,webpdf,contemporary,large]{oup-authoring-template}

\usepackage{graphicx}

% Version 1.5 of the class asks for a society's logo on the last page and
% stops the build when there is none; this says there is none.
\providecommand{\societylogo}{}

\begin{document}

\journaltitle{The Journal}
\DOI{DOI HERE}
\copyrightyear{2027}
\pubyear{2027}
\appnotes{Paper}

\firstpage{1}

\title[A Working Title]{A Working Title}

\author[1]{Your Name}
\address[1]{\orgname{Your Institution}, \orgaddress{\state{City}, \country{Country}}}
\corresp[$\ast$]{Corresponding author. \href{mailto:you@example.org}{you@example.org}}

\abstract{One paragraph saying what the work is, what you did, and what you found.
Write it last.}
\keywords{first keyword, second keyword}

\maketitle

\section{Introduction}
\label{sec:introduction}

Start here.  A citation looks like this~\cite{knuth1984}, and a reference to
another section looks like Section~\ref{sec:method}.

\section{Method}
\label{sec:method}

A numbered equation:
\begin{equation}
  \label{eq:example}
  E = mc^2.
\end{equation}

\section{Conclusion}

What the results mean, and what they do not.

\bibliographystyle{abbrvnat}
\bibliography{references}

\end{document}
